\uuid{CrPH}
\titre{Pgcd}
\auteur{bodin ; croizat ; sacre}
\organisation{exo7}
\licence{CC BY-NC-SA 4.0}
\datecreate{2022-03-08}
\niveau{L1}
\module{Algèbre}
\chapitre{Arithmétique dans Z}
\sousChapitre{Pgcd, ppcm, algorithme d'Euclide}
\difficulte{3}
\theme{pgcd}
\exo7qcmid{hilisit.arithmetique.013}

\contenu{
\begin{question}
\qclassification{103.03}
    Soit $a>0$ un entier impair qui est un multiple de $3$. Quelles sont alors les affirmations vraies ?
    \begin{answers}
        \bad{$a$ est un multiple de $6$.}
        \bad{L'écriture décimale de $a$ finit nécessairement soit par $7$ soit par $9$.}
        \good{$\pgcd(a,3) = 3$.}
        \good{$\ppcm(a,3) = a$.}
    \end{answers}
    \begin{explanations}
    Les multiples positifs de $3$ s'écrivent $3N$, avec $N$ un entier positif. Si $N=2k$ est pair, alors $3N = 6k$ est un entier pair. Donc $a$, notre entier impair multiple de $3$, s'écrit $a = 3N$ avec $N = 2k+1$ un nombre impair ; ou encore $a = 3(2k+1) = 6k + 3$.\\
    Le reste de la division euclidienne de $a$ par $6$ est $3$. Donc $a$ n'est pas un multiple de $6$.\\
    Pour $k=2$ par exemple, on a $a = 6 \times 2 + 3 = 15$ qui est un entier impair, multiple de $3$, dont l'écriture décimale ne finit ni par $7$ ni par $9$.\\
    La liste des diviseurs de $3$ se réduit à $1$ et $3$. Puisque $3$ divise $a$, $3$ est un multiple commun à $3$ et $a$ : on a donc $\pgcd(a,3) = 3$. Par conséquent, on a aussi $\ppcm(a,3) = a$ de sorte que $\pgcd(a,3) \times \ppcm(a,3) = a \times 3$.
    \end{explanations}
\end{question}
}
